\documentclass{article}
\usepackage{iclr2026_conference,times}
\usepackage{graphicx}
\usepackage{booktabs}
\usepackage{amsmath}
\usepackage{amssymb}
\usepackage{hyperref}
\usepackage{url}

\iclrfinalcopy

\title{Reward-Free World-Model Pretraining for \\ Transferable Control Representations}

\author{Ricky Bao \\
Department of Computer Science \\
University of Chicago \\
CMSC 35401: Active Representation Learning \\
\texttt{rickybao@uchicago.edu}}

\newcommand{\fig}[1]{figures/#1}

\begin{document}
\maketitle
\lhead{}
\renewcommand{\headrulewidth}{0pt}

\begin{abstract}
A world model trained without any task reward should, in principle, learn
dynamics that transfer to many downstream tasks. Whether it actually does, and
which property of the learned representation tracks that transfer,
remains poorly characterized. We study reward-free pretraining of a DreamerV3
world model on the DeepMind Control proprioceptive Walker, comparing four
pretraining strategies (none, random-policy data, Plan2Explore
latent-disagreement, and APT particle-entropy) and adapting each to three
downstream tasks (\texttt{stand}, \texttt{walk}, \texttt{run}) under a
frozen-readout protocol. Across two seeds per condition we find that
exploratory pretraining (Plan2Explore, APT) speeds early adaptation by roughly
$2$--$3\times$ over training from scratch, whereas random-policy pretraining
gives no benefit and, when frozen, actively hurts. Adaptation speed correlates
far more strongly with the coverage of the pretraining data and with
linear probes of derived dynamical quantities (time-to-fall) than with
current-state decodability, which is saturated and non-discriminative. A
time-resolved analysis shows that Plan2Explore's coverage is built
continuously by the learned exploratory policy, accumulating well after
its raw disagreement reward has collapsed, which argues against a sustained,
high-magnitude raw-disagreement reward as the driver. A minimal
protocol-calibration experiment (Plan2Explore vs.\ no-pretraining on
\texttt{walk}) shows the early-adaptation ordering carries over from
frozen-readout to standard full-adaptation, while a frozen-readout-specific
final-performance deficit disappears. The benefit is one of sample
efficiency, not a higher performance ceiling. We frame our conclusions around
what the small two-seed design can and cannot support, and are explicit that our
coverage measure is nearly collinear with ``the policy explores at all,'' so it
cannot by itself establish coverage as the causal driver of transfer.
\end{abstract}

\section{Introduction}

Model-based reinforcement learning agents such as DreamerV3
\citep{hafner2023dreamerv3} learn a latent world model and train a policy
entirely inside imagined rollouts. The world model is the agent's
representation of its environment, and a natural question is whether that
representation can be learned before any task is specified, then reused
across tasks. This is the promise of unsupervised, or reward-free,
reinforcement learning \citep{laskin2021urlb}: spend an exploration budget
learning task-agnostic dynamics, and amortize it over many downstream
objectives that share the same body and physics.

Two questions stand in the way of using this idea confidently. First, it is not
obvious that reward-free pretraining helps at all in the model-based setting: a
world model trained on a poor data distribution may model the wrong part of the
state space, and a frozen representation may simply cap downstream performance.
Second, even when pretraining helps, the field lacks a clear account of
which property of the learned representation is responsible. Exploration
methods are usually compared by their downstream return, not by the
representations they induce, so it is rarely clear whether a method helps
because it produces an accurate forward model, a linearly decodable state, or
simply broad data coverage.

We take an empirical, mechanism-focused position on a deliberately small and
controllable domain. Transfer here means reuse across reward functions on a
single Walker body and proprioceptive observation space, not
cross-morphology or cross-domain transfer. Because the 500K reward-free phase is
a sunk cost paid once, any benefit is a per-task sample-efficiency gain that only
amortizes across several downstream tasks, a point we make precise in
Section~6. We pretrain a DreamerV3 world model without task reward under four
strategies, adapt each to three Walker tasks, and pair the adaptation curves with
a battery of representation probes and a direct measure of pretraining-data
coverage. Our contributions are:

\begin{itemize}
\item A controlled four-condition comparison (no-pretrain, random, Plan2Explore,
APT) on a shared body, isolating the effect of the pretraining
exploration objective from the model and the downstream task.
\item Evidence that exploratory pretraining improves early-adaptation sample
efficiency by $2$--$3\times$, that random pretraining does not, and that under a
frozen representation random pretraining underperforms training from scratch.
\item A representation analysis showing that adaptation speed tracks data
coverage and derived-quantity probes, not current-state decodability, and a
time-resolved coverage analysis that argues against a sustained raw-disagreement
reward as the explanation for Plan2Explore's benefit.
\item A protocol-calibration experiment separating ``representation as fixed
features'' from ``representation as initialization,'' which reconciles an
apparent final-performance deficit of frozen pretraining.
\end{itemize}

Given the two-seed design and six-point correlations, these results should be
read as descriptive evidence that motivates follow-up experiments rather than as
significance claims.

\section{Background and Related Work}

\paragraph{World models and DreamerV3.}
The Dreamer family \citep{hafner2020dreamer, hafner2021dreamerv2,
hafner2023dreamerv3} learns a Recurrent State-Space Model (RSSM)
\citep{hafner2019planet} that encodes observations into a latent with a
deterministic recurrent part and a stochastic categorical part, and trains an
actor and critic on imagined latent rollouts. The world model is optimized for
reconstruction and latent prediction, so it is in principle reusable across
tasks that share dynamics; only the reward predictor and behavior are
task-specific. We use DreamerV3 because its objective separates the
task-agnostic world model from task-specific heads, matching the separation
required by a frozen-readout transfer study.

\paragraph{Intrinsic-motivation exploration.}
When no task reward is available, exploration must be driven by an intrinsic
signal. Two families are relevant here. Disagreement methods reward an agent for
visiting states where an ensemble of forward predictors disagrees, a proxy for
epistemic (reducible) uncertainty \citep{pathak2019disagreement}; Plan2Explore
\citep{sekar2020plan2explore} brings this idea into the Dreamer setting by
training the disagreement ensemble on latent features and optimizing the
exploration policy in imagination. Entropy methods instead reward state-space
coverage directly; APT \citep{liu2021apt} maximizes a particle-based
$k$-nearest-neighbor estimate of latent entropy \citep{singh2003knn}.
Related signals include prediction-error curiosity \citep{pathak2017icm} and
random network distillation \citep{burda2019rnd}. A separate URLB family pursues
behavioral diversity or successor features rather than novelty per se (e.g.,
DIAYN, \citealp{eysenbach2019diayn}; APS, \citealp{liu2021aps}); we restrict
attention to Plan2Explore and APT because they instantiate the
uncertainty-versus-coverage contrast within a single world-model agent
with the architecture held fixed, which is what our representation comparison
requires. Plan2Explore and APT are thus the natural contrast for our central
question: does \emph{uncertainty-driven} exploration induce a different, more
transferable representation than \emph{coverage-driven} exploration?

\paragraph{Unsupervised RL and adaptation protocols.}
The Unsupervised Reinforcement Learning Benchmark \citep{laskin2021urlb}
standardizes the two-stage recipe we follow: a long reward-free pretraining
phase followed by a short reward-driven adaptation phase, on shared-body domains
including Walker. URLB adapts by continuing to train the whole agent
(\emph{full-adaptation}). We instead make \emph{frozen-readout} primary (the
encoder and RSSM are frozen and only a fresh task head is trained) because it
measures the representation as a fixed feature space rather than as an
initialization, and then calibrate the two protocols against each other. A
recurring URLB finding is that representation-centric methods can rank
differently under the two protocols, which our calibration revisits directly.

\paragraph{Probing representations.}
Linear probes \citep{alain2016probes} are a standard tool for asking what a
frozen representation encodes. In a proprioceptive control setting the
observation is already close to the simulator state, so probing the
current state is near-trivial; we therefore emphasize future-state
probes, derived dynamical quantities, and reward decodability, which are not
linear read-outs of a single frame.

\section{Research Questions and Approach}

We address three questions:
\textbf{(RQ1)} Does reward-free world-model pretraining improve downstream
adaptation, and does the exploration objective matter?
\textbf{(RQ2)} Which property of the pretrained representation is most
associated with adaptation speed: forward-predictive accuracy, decodability
of state and derived quantities, reward decodability, or data coverage?
\textbf{(RQ3)} Do uncertainty-driven (Plan2Explore) and coverage-driven (APT)
exploration induce measurably different transfer?

\paragraph{Conditions.}
All conditions share the DreamerV3 model and the Walker body, and differ only in
how the world model is pretrained for 500K reward-free environment steps:
\textbf{C1 (no-pretrain)} trains DreamerV3 from scratch on the downstream task;
\textbf{C2 (random)} trains the world model on data from a uniform-random
policy, with no actor; \textbf{C3 (Plan2Explore)} adds an ensemble of one-step
latent predictors and trains an exploration actor on their disagreement;
\textbf{C4 (APT)} trains an exploration actor on the $k$-NN particle-entropy of
the latent. C1 is the no-pretraining baseline (not a strict lower bound, since
at the adaptation budget it can beat a frozen pretrained model on final return);
C2 isolates ``a world model trained on passive data'' from the effect of active
exploration.

\paragraph{Intrinsic rewards (implementation).}
For C3 the ensemble predicts the next \emph{deterministic posterior feature}
(the recurrent state concatenated with the categorical posterior probabilities)
rather than the sampled one-hot latent. We found the sampled-latent target
degenerate: its irreducible sampling noise floors the ensemble loss while the
optimal predictor collapses to the posterior mean, so disagreement vanishes even
though the loss does not. Predicting the deterministic feature removes that
aleatoric noise, leaving epistemic disagreement. Members use independent
initialization and bootstrap masking. For C4 the reward at each imagined latent
is $\log(c + \tfrac{1}{k}\sum_{j \in \mathrm{kNN}} \lVert z - z_j \rVert)$
computed over the imagination batch. Both intrinsic rewards are stop-gradient'd
and consumed in place of the task reward head during imagination. Both are
\emph{adapted} variants of the originals: canonical Plan2Explore predicts the
next encoder embedding and APT estimates entropy in a separate encoder space,
whereas we predict the RSSM posterior feature and estimate entropy in the
world-model latent. These choices fit the proprioceptive, latent-only setting,
but cross-paper comparisons of absolute numbers should be read with the
adaptation in mind.

\paragraph{Adaptation protocol.}
Frozen-readout is primary: we load the pretrained encoder, RSSM, and decoder,
freeze them, and train a fresh actor, critic, and reward predictor on the
downstream reward, using DreamerV3's imagination through the frozen RSSM. The
adaptation budget is $A=250$K steps. The \textbf{primary metric} is
early-adaptation AUC: the mean episode return over the first $125$K steps (the
first half of $A$), equivalently a normalized area under the return curve, which
reflects both how high adaptation starts and how fast it climbs. The
\textbf{secondary metric} is final return near $A$. We report both because they
answer different questions: how fast versus how high. With only two seeds per
condition we report per-seed values and seed spread rather than significance
tests, and make no significance claims, in keeping with small-sample reporting
guidance for deep RL \citep{agarwal2021rliable}.

\paragraph{Representation analysis.}
We collect one common held-out trajectory set from a fixed uniform-random policy
with logged MuJoCo physics, and encode it through each condition's frozen
world model, so all conditions are probed on identical underlying states. On
these features we fit ridge probes for current and future simulator state, gait
phase, time-to-fall, and per-task reward, reporting held-out $R^2$. Separately,
we measure coverage of each condition's pretraining replay as a $k$-NN
particle entropy of the proprioceptive observations (an almost-complete function
of the simulator state in this domain), on a shared standardizer so conditions
are comparable. To probe mechanism we also bucket each replay buffer by
collection time and recompute coverage per window.

\section{Experimental Setup}

\paragraph{Environment and model.}
We use the DeepMind Control Suite \citep{tassa2018dmc} Walker in
proprioceptive mode (no pixels; observation = orientations, torso height,
velocities), simulated in MuJoCo \citep{todorov2012mujoco}. The agent is the
public DreamerV3 (JAX) at its smallest preset (\texttt{size1m}: RSSM
deterministic state $512$, $32$ categorical latents of $4$ classes,
$64$-unit MLPs), trained at \texttt{train\_ratio} $1024$. Pretraining is 500K
reward-free steps; adaptation is 250K steps; C1 from-scratch runs to 500K, which
also provides a full-performance reference.

\paragraph{Conditions, tasks, seeds.}
Four conditions $\times$ three tasks, two seeds each. C1/\texttt{walk} reuses two
existing 500K from-scratch Walker-walk runs, so no new C1/\texttt{walk} job is
needed. This gives 6 reward-free pretraining runs, 4 new C1 runs, and 18
frozen-readout adaptation runs. All metrics are reported with the per-seed
spread, never as single point estimates.

\paragraph{Protocol calibration.}
To test whether frozen-readout conclusions transfer to the standard URLB
protocol, we additionally run C3 (Plan2Explore) full-adaptation on
\texttt{walk} for two seeds, in which the world model continues to train during
adaptation. This is the minimal version of the calibration; a full four-way
calibration is left to future work.

\section{Results}

\subsection{RQ1: Exploratory pretraining speeds adaptation; random pretraining
does not}

Table~\ref{tab:auc} reports the primary metric. Exploratory pretraining (C3,
C4) improves early-adaptation AUC over training from scratch by roughly
$2$--$3\times$ on every task: on \texttt{stand}, C1 reaches $321$ versus $676$
(C3) and $740$ (C4); on \texttt{run}, $52$ versus $141$ and $167$. Random
pretraining (C2) is essentially identical to C1 ($\approx 312$ vs.\ $321$ on
\texttt{stand}). We state these as multiplicative ranges over two-seed means and
report the per-seed values in Appendix Table~\ref{tab:perseed}; with two seeds we
claim a consistent direction and rough magnitude, not statistical significance.
Figure~\ref{fig:auc} shows the same picture per task, and
Figure~\ref{fig:curves} shows the underlying learning curves: the exploratory
conditions start higher and climb faster.

\begin{table}[t]
\centering
\caption{Early-adaptation AUC (mean episode return over the first 125K steps;
mean $\pm$ half the seed range, $n=2$). Higher is faster adaptation.}
\label{tab:auc}
\begin{tabular}{lccc}
\toprule
Condition & \texttt{stand} & \texttt{walk} & \texttt{run} \\
\midrule
C1 no-pretrain   & $321 \pm 58$ & $180 \pm 59$ & $52 \pm 3$ \\
C2 random        & $312 \pm 49$ & $175 \pm 6$  & $63 \pm 13$ \\
C3 Plan2Explore  & $676 \pm 39$ & $\mathbf{381 \pm 39}$ & $141 \pm 3$ \\
C4 APT           & $\mathbf{740 \pm 36}$ & $323 \pm 59$ & $\mathbf{167 \pm 19}$ \\
\bottomrule
\end{tabular}
\end{table}

\begin{figure}[t]
\centering
\includegraphics[width=0.78\textwidth]{\fig{01_early_auc_by_task.png}}
\caption{Early-adaptation AUC by condition and task. Plan2Explore and APT
improve over both C1 (no-pretrain) and C2 (random); random pretraining matches
the from-scratch baseline. Bars show the two-seed mean; markers show individual
seeds.}
\label{fig:auc}
\end{figure}

The secondary metric tells a complementary and important story
(Table~\ref{tab:final}). By $240$K steps, from-scratch C1 has largely caught up:
on \texttt{walk} it reaches $754$, ahead of frozen C3 ($535$) and C4 ($765$,
essentially tied). The pretraining benefit is therefore primarily one of
sample efficiency rather than asymptotic performance at this budget. Two
findings stand out. First, frozen C2 (random) collapses far below C1 on every
task ($144$ vs.\ $754$ on \texttt{walk}); its early and final values are nearly
equal, meaning it plateaus almost immediately. A frozen world model trained on
random-policy data, which mostly captures a flailing or fallen Walker, is a
worse fixed representation than learning from scratch. Second, frozen C3 trails
C1 at convergence on \texttt{walk} and \texttt{run} even though it adapts faster
early, a tension we resolve in Section~\ref{sec:calib}.

\begin{table}[t]
\centering
\caption{Final return (mean episode return over steps 220--240K; mean $\pm$ half
the seed range). C1 is measured at the same step budget.}
\label{tab:final}
\begin{tabular}{lccc}
\toprule
Condition & \texttt{stand} & \texttt{walk} & \texttt{run} \\
\midrule
C1 no-pretrain   & $877 \pm 36$ & $754 \pm 19$ & $246 \pm 54$ \\
C2 random        & $325 \pm 48$ & $144 \pm 4$  & $90 \pm 35$ \\
C3 Plan2Explore  & $887 \pm 69$ & $535 \pm 89$ & $176 \pm 31$ \\
C4 APT           & $929 \pm 19$ & $765 \pm 78$ & $243 \pm 40$ \\
\bottomrule
\end{tabular}
\end{table}

\subsection{RQ2: Coverage and derived-quantity probes track transfer; current
state does not}

Table~\ref{tab:repr} summarizes the representation analysis on the common
held-out set, and Figure~\ref{fig:probe} visualizes it. The current-state probe
is saturated: every condition decodes the simulator state with $R^2 \approx
0.94$, and random pretraining is in fact marginally highest ($0.953$). Current
linear decodability therefore carries essentially no signal about transfer in
this proprioceptive domain. The strong negative correlation between the
current-state probe and AUC ($r\approx -0.88$) is an artifact of this saturation
(all values near a ceiling, the slowest condition marginally above it) and
should not be read as a real effect.

The discriminative metrics are data coverage and the derived-quantity probes.
Coverage (particle entropy of the pretraining replay) orders the conditions
random $1.10 <$ APT $1.26 <$ Plan2Explore $1.39$, and the time-to-fall probe
$R^2$ orders them similarly ($0.55 < 0.63 \le 0.65$). Across the six pretrained
(condition, seed) points, early-adaptation AUC correlates $+0.68$ to $+0.90$
with coverage and $+0.67$ to $+0.90$ with time-to-fall (Table~\ref{tab:repr}).
Three caveats apply here. First, with $n=6$ these correlations are descriptive,
not significant; we read their sign and rough size, not their exact value.
Second, and more fundamentally, our coverage measure is nearly collinear with
``the policy explores at all'': random barely moves and the two exploratory
conditions both move, so any metric separating active from passive policies will
correlate with AUC. With three conditions the correlation cannot isolate
coverage as the \emph{causal} driver; it largely restates the
exploratory-versus-passive split the conditions were built to create. Third, the
$k{=}20$ future-state probe appears with a high positive correlation in
Table~\ref{tab:repr}, but its absolute $R^2$ is within noise of zero for every
condition, so that correlation rides on the sign of negligible numbers and is no
more interpretable than the saturated current-state probe.

\begin{table}[t]
\centering
\caption{Representation metrics on the common held-out set (condition means
over 2 seeds), and their correlation with early-adaptation AUC across the six
pretrained (condition, seed) points (range over the three tasks, $n=6$). C1 has
no pretrained world model and is omitted.}
\label{tab:repr}
\begin{tabular}{lcccc}
\toprule
Metric & C2 random & C3 P2E & C4 APT & corr.\ with AUC \\
\midrule
Current-state probe $R^2$        & 0.953 & 0.941 & 0.944 & $-0.85$ to $-0.90$ (artifact) \\
Future-state probe $R^2$ ($k{=}20$) & $-0.003$ & 0.002 & 0.001 & $+0.84$ to $+0.92$ \\
Time-to-fall probe $R^2$         & 0.554 & 0.627 & 0.647 & $+0.67$ to $+0.90$ \\
Reward probe $R^2$ (\texttt{walk}) & 0.615 & 0.651 & 0.641 & $+0.48$ to $+0.62$ \\
Coverage (particle entropy)      & 1.100 & \textbf{1.386} & 1.262 & $+0.68$ to $+0.86$ \\
\bottomrule
\end{tabular}
\end{table}

\begin{figure}[t]
\centering
\includegraphics[width=0.92\textwidth]{\fig{04_probe_and_coverage_summary.png}}
\caption{Representation summary. The current-state probe is saturated across
conditions (left), while data coverage clearly separates the methods (right).
Adaptation speed tracks the latter, not the former.}
\label{fig:probe}
\end{figure}

\subsection{RQ3: Plan2Explore and APT transfer comparably, and coverage
outlasts the disagreement reward}

On adaptation, C3 and C4 are broadly comparable: APT leads on \texttt{stand} and
\texttt{run}, Plan2Explore on \texttt{walk}, and their seed ranges overlap on
\texttt{stand} and \texttt{walk}. We therefore cannot separate the two methods'
downstream effect at two seeds, which is itself the answer to RQ3 at this
budget: uncertainty-driven and coverage-driven exploration produce similar
transfer here. The uncertainty-driven method (Plan2Explore) in fact produced
the highest data coverage, more than the coverage-driven method (APT),
so coverage and the nominal objective do not line up in the naive way.

The timing analysis sharpens the mechanism. Plan2Explore's raw disagreement
reward collapses by roughly $100\times$ early in pretraining (from $\sim
0.046$ to $\sim 0.0005$ by 250K steps) as the world model reduces its
uncertainty. Yet its replay coverage keeps rising monotonically through the end
of pretraining (Figure~\ref{fig:timing}, Figure~\ref{fig:covtime}), long after
the reward has vanished, whereas random-policy coverage stays flat throughout.
The \emph{learned exploratory policy} thus generates this coverage continuously,
well after the intrinsic reward has decayed, rather than in an early burst that
is later retained. This argues against a sustained, high-magnitude
\emph{raw}-disagreement reward as the driver, though we measured the raw signal,
not the percentile-normalized advantage the actor actually optimizes, so a
normalized or relative-uncertainty effect is not excluded. The coverage account
is consistent with the data, but does not prove it,
and we give equal weight to a competing explanation: the exploratory policy
becomes an active, upright mover, so its world model simply sees the
task-relevant locomotion regime that the random policy never reaches. Under this
view the operative variable is not breadth of coverage per se but whether
pretraining reaches the right region of state space at all. Our data cannot
separate ``broad coverage'' from ``reached the task-relevant regime,'' because
the conditions that have one have the other; distinguishing them (Section~6) is
the natural next step.

\begin{figure}[t]
\centering
\includegraphics[width=0.82\textwidth]{\fig{05_p2e_timing_coverage_vs_disagreement.png}}
\caption{Plan2Explore timing. The raw disagreement reward collapses early in
pretraining, while replay coverage continues to grow, indicating the useful data
distribution is sustained by the learned policy rather than by the reward
signal.}
\label{fig:timing}
\end{figure}

\begin{figure}[t]
\centering
\includegraphics[width=0.72\textwidth]{\fig{03_coverage_over_time_mean.png}}
\caption{Coverage over pretraining time, bucketed by collection step. Random
coverage is flat; Plan2Explore and APT coverage rise monotonically and peak in
the final window.}
\label{fig:covtime}
\end{figure}

\subsection{Protocol calibration: the early-adaptation gap survives
full-adaptation; the final-return deficit does not}
\label{sec:calib}

Table~\ref{tab:calib} compares C1 and C3 on \texttt{walk} under both protocols.
On the primary metric the two protocols agree for this C1-vs-C3 contrast:
C3 beats C1 early under frozen-readout ($+200$ AUC) and by even more under
full-adaptation ($+382$), so the early-adaptation gap is not a frozen-readout
artifact for the one contrast we calibrated. We do not claim
protocol-invariance beyond C1-vs-C3 on \texttt{walk}. On final return the
ordering flips: under freezing C3 trails C1 by $219$, but under
full-adaptation C3 leads by $191$ and reaches near-ceiling ($945$). The
frozen-readout final deficit reflects the cost of holding the world model fixed,
not a deficiency of the pretrained dynamics. We did not run C2
full-adaptation, so we do not extend this explanation to C2's collapse, which
under freezing could equally reflect genuinely poor random-policy dynamics. Plan2Explore's world model is a
strong \emph{initialization} even where it is a mediocre \emph{fixed feature
space}, precisely the ``representation as initialization vs.\ fixed features''
distinction the calibration was designed to expose. Figure~\ref{fig:calib}
visualizes the three curves.

\begin{table}[t]
\centering
\caption{Protocol calibration on \texttt{walk} (2 seeds). For the C1-vs-C3
contrast, the early-AUC ordering matches across protocols; the final-return
ordering does not.}
\label{tab:calib}
\begin{tabular}{lcc}
\toprule
Setting & Early-AUC & Final return \\
\midrule
C1 from-scratch (inherently full) & 180 & 754 \\
C3 frozen-readout                  & 381 & 535 \\
C3 full-adaptation                 & \textbf{562} & \textbf{945} \\
\midrule
C1$\to$C3 gap (frozen)             & $+200$ & $-219$ \\
C1$\to$C3 gap (full)               & $+382$ & $+191$ \\
\bottomrule
\end{tabular}
\end{table}

\begin{figure}[t]
\centering
\includegraphics[width=0.62\textwidth]{\fig{06_fulladapt_walk_calibration.png}}
\caption{Walk adaptation under both protocols. Full-adaptation C3 dominates both
frozen C3 and from-scratch C1 on speed and final return.}
\label{fig:calib}
\end{figure}

\section{Discussion and Conclusion}

We asked whether reward-free world-model pretraining yields transferable control
representations, which representational property tracks the transfer, and
whether the exploration objective matters. On a controlled DMC Walker study with
four pretraining conditions, three tasks, and matched probes, we find that
exploratory pretraining (Plan2Explore, APT) improves early-adaptation sample
efficiency by $2$--$3\times$, that random-policy pretraining does not and is
actively harmful as a frozen representation, and that adaptation speed tracks
the coverage of the pretraining data and derived-quantity probes rather than the
saturated current-state probe (RQ1, RQ2). Uncertainty-driven and coverage-driven
exploration transfer comparably (RQ3), and a time-resolved analysis argues
against a sustained, high-magnitude raw-disagreement reward as the driver for
Plan2Explore: its coverage is built continuously by the learned policy, even
after its raw disagreement reward collapses. A minimal calibration (C1 vs.\ C3
on \texttt{walk}) shows that the early-adaptation ordering carries over to
standard full-adaptation, while a frozen-readout-specific final-performance
deficit disappears.

\paragraph{Implications.}
The results suggest a few lessons beyond this Walker setup. First, in proprioceptive control the
representation property that tracks transfer is not the linear decodability of
the current state, which saturates, but the data distribution the exploration
policy induces (which here we cannot cleanly separate from simply reaching the
task-relevant regime). Either way, this argues for evaluating
exploration methods by the distribution they collect, not only by downstream
return. Second, the choice of adaptation protocol is itself a finding, not a
nuisance: frozen-readout is a conservative stress test that can invert
final-performance orderings relative to full-adaptation, so transfer claims
should state which protocol they assume. Third, because the gain is sample
efficiency rather than a higher ceiling, a 500K reward-free pretraining phase
only ``pays for itself'' when amortized across several downstream tasks: it is a
one-time cost that buys a faster start on each task, and at our budgets the early
saving is on the order of $10^5$ steps per task, so the economics favor
pretraining only when the number of downstream tasks sharing the body is large
enough to outweigh the pretraining steps.

\paragraph{Limitations.}
The study uses two seeds per condition and six points per correlation, so all
correlations are descriptive rather than significant and the C3-vs-C4 ordering is
unresolved; the seed spread we report makes this visible. The
coverage-correlation is partly confounded by the binary exploratory-vs-passive
split, and because our coverage measure is nearly collinear with
``the policy is an active explorer,'' the analysis cannot attribute transfer to
coverage as opposed to simply reaching the task-relevant regime.
Coverage-mediation of transfer therefore remains a hypothesis: the clean causal
test (adapting from intermediate pretraining checkpoints) requires
checkpoint retention we did not enable and is left to future work. The protocol
calibration is minimal (Plan2Explore on \texttt{walk} only); a full four-way,
three-task calibration would be needed to claim protocol-invariance broadly. The
pretraining budget (500K) is short relative to URLB's $\sim$2M; this especially
weakens the negative C2 result, since a random policy might collect more useful
data given much longer, so we read ``random pretraining does not help'' as
budget-conditional rather than absolute. Finally, the domain is a single shared
body, and with two seeds we report direction and spread rather than
significance, so absolute magnitudes and the unresolved C3-versus-C4 ordering
should not be over-interpreted.

\paragraph{Future work.}
The most informative next experiment is adaptation from time-stamped pretraining
checkpoints, which would convert the coverage-timing observation into a causal
test. Extending the four-way calibration and adding seeds would resolve C3
versus C4. Finally, disentangling ``broad coverage'' from ``task-relevant active
locomotion'' (for example by measuring coverage restricted to upright, moving
states) would clarify whether exploration helps by breadth or by reaching the
right regime.

\subsubsection*{Reproducibility}
Code for world-model pretraining, adaptation, probing, coverage, and synthesis
is available at
\url{https://github.com/shalalababa/dreamerv3/tree/arl-exploration}. The branch
contains the scripts and report figures used here; raw run outputs are omitted
from the repository. Reported numbers are computed by the probing pipeline
(\texttt{probing/}) from logged run metrics; tables cite condition means with
per-seed spread.

\bibliographystyle{iclr2026_conference}
\bibliography{final_report}

\appendix
\section{Per-seed results}
Table~\ref{tab:perseed} gives the per-seed early-adaptation AUC behind the
condition means in Table~\ref{tab:auc}, showing the two-seed spread directly.

\begin{table}[h]
\centering
\caption{Per-seed early-adaptation AUC.}
\label{tab:perseed}
\begin{tabular}{llccc}
\toprule
Condition & Seed & \texttt{stand} & \texttt{walk} & \texttt{run} \\
\midrule
C1 & 1 & 378.8 & 120.9 & 54.4 \\
C1 & 2 & 263.3 & 239.2 & 49.1 \\
C2 random & 1 & 361.3 & 181.0 & 76.5 \\
C2 random & 2 & 262.8 & 169.7 & 50.0 \\
C3 P2E & 1 & 714.7 & 341.6 & 144.2 \\
C3 P2E & 2 & 637.1 & 419.5 & 137.2 \\
C4 APT & 1 & 775.8 & 264.2 & 185.9 \\
C4 APT & 2 & 703.5 & 381.4 & 148.7 \\
\bottomrule
\end{tabular}
\end{table}

\begin{table}[h]
\centering
\caption{Per-seed final return (mean episode return over steps 220--240K),
behind the condition means in Table~\ref{tab:final}.}
\label{tab:perseed-final}
\begin{tabular}{llccc}
\toprule
Condition & Seed & \texttt{stand} & \texttt{walk} & \texttt{run} \\
\midrule
C1 & 1 & 913 & 735 & 300 \\
C1 & 2 & 841 & 773 & 192 \\
C2 random & 1 & 372 & 148 & 125 \\
C2 random & 2 & 277 & 140 & 55 \\
C3 P2E & 1 & 955 & 624 & 207 \\
C3 P2E & 2 & 818 & 446 & 145 \\
C4 APT & 1 & 948 & 687 & 282 \\
C4 APT & 2 & 910 & 844 & 203 \\
\bottomrule
\end{tabular}
\end{table}

\begin{table}[h]
\centering
\caption{Per-seed representation metrics on the common held-out set (probe
$R^2$ and replay coverage), behind the condition means in Table~\ref{tab:repr}.
C1 has no pretrained world model. The future-state probe ($k{=}20$) is omitted
because its $R^2$ is within noise of zero for every condition and seed (cf.\
Table~\ref{tab:repr}) and is therefore not interpretable per seed.}
\label{tab:perseed-repr}
\small
\begin{tabular}{llcccc}
\toprule
Cond. & Seed & state $R^2_{k=0}$ & TTF $R^2$ & rew$_{\mathrm{walk}}$ $R^2$ & cov. \\
\midrule
C2 random & 1 & 0.953 & 0.578 & 0.611 & 1.099 \\
C2 random & 2 & 0.953 & 0.530 & 0.620 & 1.102 \\
C3 P2E & 1 & 0.942 & 0.658 & 0.676 & 1.433 \\
C3 P2E & 2 & 0.941 & 0.597 & 0.626 & 1.340 \\
C4 APT & 1 & 0.944 & 0.632 & 0.627 & 1.227 \\
C4 APT & 2 & 0.945 & 0.663 & 0.656 & 1.296 \\
\bottomrule
\end{tabular}
\end{table}

\paragraph{Full-adaptation calibration, per seed.} C3 (Plan2Explore)
full-adaptation on \texttt{walk}: seed~1 early-AUC $444$, final $943$; seed~2
early-AUC $679$, final $948$ (means $562$ / $945$ in Table~\ref{tab:calib}).

\section{Hyperparameters and reproducibility details}
\begin{itemize}
\item \textbf{Environment.} DMC Walker, proprioceptive observations
(orientations, torso height, velocities; $24$-d), action repeat $1$, episode
length $1000$, MuJoCo physics.
\item \textbf{Agent.} DreamerV3 \texttt{size1m}: RSSM deterministic state $512$,
$32$ stochastic categoricals of $4$ classes, encoder/decoder MLPs $64$ units
$\times\,3$ layers, symlog inputs. Batch $16 \times 64$; optimizer learning rate
$4\times10^{-5}$; \texttt{train\_ratio} $1024$; $16$ parallel envs; imagination
horizon $15$; discount horizon $333$.
\item \textbf{Budgets / seeds.} Pretraining $500$K reward-free steps; adaptation
$A=250$K; C1 from-scratch to $500$K. Two seeds per condition, RNG seeds
$\{1,2\}$.
\item \textbf{Plan2Explore (C3).} Ensemble of $8$ one-step predictors ($2$
layers $\times\,256$ units), target = next deterministic posterior feature
(deter $\oplus$ softmax-posterior, $640$-d), bootstrap-mask probability $0.8$,
intrinsic-reward scale $10^3$ before percentile normalization.
\item \textbf{APT (C4).} $k$-NN particle entropy over the $640$-d world-model
latent with $k=12$, $c=1.0$, computed within each imagination batch.
\item \textbf{Frozen-readout.} Encoder + RSSM + decoder loaded from the
pretraining checkpoint (parameter regex \texttt{\^{}(enc|dyn|dec)/}) and frozen;
a fresh actor, critic, and reward head are trained.
\item \textbf{Probes.} Common held-out set: $40$ episodes from a uniform-random
policy with logged physics. Ridge probes with $\lambda \in \{10^{-3},10^{-1},
1,10,10^{3}\}$ selected on a $20\%$ validation split; test fraction $0.3$; sites
encoder / posterior / prior / open-loop imagination at $k\in\{1,5,20\}$; target
horizons $\{0,1,5,20\}$.
\item \textbf{Coverage.} Particle entropy with $k=12$, $c=1.0$, up to $3000$
frames per buffer; $2$-D histogram entropy over the top two principal components
($24$ bins); a single standardizer and PCA fit on the pooled frames so
conditions share a scale; $6$ equal time buckets for the coverage-over-time
curve.
\item \textbf{Metrics.} Episode return is logged at each episode completion.
Early-adaptation AUC is the trapezoidal integral of (step, episode-return) over
the first $125$K steps divided by the integrated step span (a normalized mean);
final return is the mean episode return over steps $220$--$240$K. Both are
computed per seed and reported as the two-seed mean with the seed half-range.
\end{itemize}

\section{Adaptation learning curves}
Figure~\ref{fig:curves} shows the full adaptation learning curves underlying the
AUC summary.

\begin{figure}[h]
\centering
\includegraphics[width=0.95\textwidth]{\fig{02_adaptation_learning_curves.png}}
\caption{Adaptation learning curves (return vs.\ environment steps) by condition
and task. Exploratory conditions start higher and rise faster; from-scratch C1
converges high by the end of the budget; frozen random plateaus low.}
\label{fig:curves}
\end{figure}

\end{document}
